Skin cancer, and melanoma in particular, is highly treatable when detected early,
but the prognosis deteriorates dramatically at advanced stages; at the same time,
access to specialist dermatological care is unevenly distributed worldwide. This
thesis presents DermAI, an end-to-end, artificial-intelligence-based skin lesion
analysis system that supports early screening through a tool available both on
smartphones and in web browsers.

The research part of the thesis trains a Vision Transformer (ViT) model on the
ISIC 2024 dataset to classify skin lesions as benign or malignant, with particular
attention to handling the dataset's severe class imbalance. The model is compared
against a ResNet-18 convolutional baseline trained under identical conditions and
a zero-shot Gemini large language model, using the official ISIC 2024 metric
complemented by statistical significance testing and calibration analysis. The
engineering part extends the trained model into a complete application, comprising
mobile and web clients, an MCP layer serving the model, Gemini-based image
validation and natural-language explanation, and authenticated, per-user cloud
storage.

\textbf{Keywords}: skin cancer, melanoma, deep learning, Vision Transformer, ISIC
2024, class imbalance, large language model, mobile health
